(1) 两电子总自旋三重态包含 \(|\uparrow,\uparrow\rangle\)、\(|\downarrow,\downarrow\rangle\)、\((|\uparrow,\downarrow\rangle+|\downarrow,\uparrow\rangle)/\sqrt2\)，单态为 \(|S\rangle=(|\uparrow,\downarrow\rangle-|\downarrow,\uparrow\rangle)/\sqrt2\)，相对符号应按固定费米子算符排列理解。双占据对称态为 \(|D_+\rangle=(|\uparrow\downarrow,0\rangle+|0,\uparrow\downarrow\rangle)/\sqrt2\)，另有反对称双占据态 \(|D_-\rangle\)。
(2) 跃迁项守恒总自旋，只把 \(|S\rangle\) 与 \(|D_+\rangle\) 耦合，取合适基矢相位后矩阵为 \(H_{S,D}=\begin{pmatrix}0&-2t\\-2t&U\end{pmatrix}\)，其本征值 \(E_\pm=[U\pm\sqrt{U^2+16t^2}]/2\)。三重态不发生双占据，且对称跃迁被交换对称性阻断，能量 \(E_T=0\)；反对称双占据态独立处于能量 \(U\)。
(3) 基态能量为 \(E_-= [U-\sqrt{U^2+16t^2}]/2<0\)。单态—三态交换能隙 \(J=E_T-E_-=[\sqrt{U^2+16t^2}-U]/2\)。若 \(U\gg t\)，展开 \(\sqrt{U^2+16t^2}=U+8t^2/U+O(t^4/U^3)\)，故 \(J=4t^2/U+O(t^4/U^3)\)，产生有效反铁磁超交换。
(4) 参数 \(U\) 耦合的算符 \(\partial H/\partial U=\sum_jn_{j\uparrow}n_{j\downarrow}\)，因此基态总双占据概率 \(P_D=\partial E_-/\partial U=\frac12[1-U/\sqrt{U^2+16t^2}]\)。\(U\to0^+\) 时 \(P_D\to1/2\)；强排斥时 \(P_D\simeq4t^2/U^2\)，双占据被抑制，三重态严格禁止同轨道反平行占据的贡献，体现费米交换及相互作用共同作用。
